\section{ماژول فشرده‌سازی مدولاسیون رتبه در ابعاد مقیاس‌یافته (\lr{SDRM Compressor})}
\label{sec:sdrm_compressor}

\subsection{معرفی و انگیزه}
این مکانیزم بر مبنای مبانی نظری پیشرفته در ادبیات توجه و فشرده‌سازی \cite{ref_sdrmcompressor_2024} توسعه یافته است.

در الگوریتم‌های سنتی فشرده‌سازی کم‌رتبه، رتبه تقریب (\lr{Rank}) معمولاً به صورت یک عدد ثابت برای تمامی لایه‌های شبکه در نظر گرفته می‌شود. این رویکرد غیرتطبیقی، تفاوت‌های ساختاری میان لایه‌ها را نادیده می‌گیرد؛ به طوری که برخی لایه‌ها با ابعاد ماتریسی کوچک به شدت فشرده شده و دچار افت شدید اطلاعات می‌شوند، در حالی که لایه‌های بسیار بزرگ با ابعاد پنهان وسیع، سهم کمی از ظرفیت فشرده‌سازی را آزاد می‌کنند.

ماژول \lr{SDRM} (\lr{Scaled Dimension Rank Modulation}) با هدف بهینه‌سازی توزیع بودجه پارامتری طراحی شده است. ایده محوری این روش، مدولاسیون پویای رتبه تجزیه متناسب با ابعاد هندسی هر لایه و بودجه پارامتری کلی است. این سازوکار تضمین می‌کند که تخصیص رتبه به صورت تحلیلی و به گونه‌ای صورت پذیرد که ضمن دستیابی به کاهش قطعی در پارامترها و محاسبات شناور (\lr{FLOPs})، توازن عملکردی مدل در تمامی لایه‌ها حفظ گردد.

\subsection{مبانی نظری و فرمول‌بندی ریاضی}
فرض کنید لایه خطی مورد نظر دارای ماتریس وزن $\mathbf{W} \in \mathbb{R}^{d_{\text{out}} \times d_{\text{in}}}$ با تعداد پارامتر اولیه $P_{\text{orig}} = d_{\text{in}} \cdot d_{\text{out}}$ باشد.

\begin{enumerate}
    \item \textbf{تخصیص بودجه و تعیین رتبه مدوله‌شده:}
    اگر نسبت بودجه مجاز برای پارامترهای لایه پس از فشرده‌سازی با $\beta \in (0, 1)$ (\lr{total\_param\_ratio}) مشخص گردد، بودجه مجاز پارامترها برابر است با:
    \begin{equation}
    P_{\text{allowed}} = \beta \cdot (d_{\text{in}} \cdot d_{\text{out}})
    \end{equation}
    در تجزیه به ساختار دو فاکتوره با رتبه $r$، تعداد پارامترهای حاصله برابر است با $P_{\text{fact}} = r \cdot (d_{\text{in}} + d_{\text{out}})$. از این رو، رتبه بهینه مدوله‌شده به صورت مستقیم از رابطه زیر تعیین می‌شود:
    \begin{equation}
    r_{\text{target}} = \left\lfloor \frac{P_{\text{allowed}}}{d_{\text{in}} + d_{\text{out}}} \right\rfloor = \left\lfloor \frac{\beta \cdot d_{\text{in}} d_{\text{out}}}{d_{\text{in}} + d_{\text{out}}} \right\rfloor
    \end{equation}

    \item \textbf{اعمال محدودیت‌های مرزی و سقف کاهش پارامتر:}
    جهت تضمین کاهش قطعی حجم و ممانعت از تنزل بیش از حد رتبه، مقدار $r$ با روابط زیر مهار می‌گردد:
    \begin{equation}
    r_{\text{cap}} = \max\left(1, \left\lfloor \frac{d_{\text{in}} d_{\text{out}}}{d_{\text{in}} + d_{\text{out}}} \right\rfloor - 1\right)
    \end{equation}
    \begin{equation}
    r = \max\left(r_{\text{min}}, \min\left(r_{\text{target}}, \min(d_{\text{out}}, d_{\text{in}}) - 1, r_{\text{cap}}\right)\right)
    \end{equation}

    \item \textbf{تجزیه ماتریسی و ساخت لایه \lr{SVDLinear}:}
    با تثبیت رتبه مدوله‌شده $r$، تجزیه مقادیر منفرد بر روی ماتریس وزن اعمال شده و دو فاکتور $\mathbf{V}_{\text{down}} \in \mathbb{R}^{r \times d_{\text{in}}}$ و $\mathbf{U}_{\text{up}} \in \mathbb{R}^{d_{\text{out}} \times r}$ استخراج می‌شوند:
    \begin{equation}
    \mathbf{W} \approx \mathbf{U}_r \mathbf{\Sigma}_r \mathbf{V}_r^\top \implies \mathbf{V}_{\text{down}} = \mathbf{\Sigma}_r^{1/2} \mathbf{V}_r^\top, \quad \mathbf{U}_{\text{up}} = \mathbf{U}_r \mathbf{\Sigma}_r^{1/2}
    \end{equation}
    انتشار خروجی با رابطه خطی زنجیره‌ای بدون نیاز به ساخت ماتریس متراکم اولیه انجام می‌شود:
    \begin{equation}
    \mathbf{y} = \mathbf{x} \mathbf{V}_{\text{down}}^\top \mathbf{U}_{\text{up}}^\top + \mathbf{b}
    \end{equation}
\end{enumerate}

\paragraph{تحلیل طیفی و نامساوی‌های پایداری خطا}
برای بررسی میزان دقت تقریب ماتریس وزن یا نگاشت توجه در ماژول \lr{sdrm\_compressor}، تفاضل ماتریس فاکتورگیری‌شده از نگاشت متراکم اصلی بر اساس نامساوی هافدینگ-فون‌نیومان ارزیابی می‌گردد. فرض کنید $\mathbf{W} \in \mathbb{R}^{d_{\text{in}} \times d_{\text{out}}}$ ماتریس وزن هدف باشد که توسط تقریب کم‌مرتبه یا پراکنده با رتبه موثر $r$ جایگزین شده است. کران بالای خطای بازسازی بر مبنای نرم فروبنیوس به صورت زیر تصریح می‌شود:
\begin{equation}
\label{eq:spectral_bound_sdrm_compressor}
\| \mathbf{W} - \mathbf{W}_{r} \|_F = \sqrt{\sum_{i=r+1}^{\min(d_{\text{in}}, d_{\text{out}})} \sigma_i^2(\mathbf{W})} \le \epsilon \| \mathbf{W} \|_F
\end{equation}
که در آن $\sigma_i(\mathbf{W})$ بیانگر $i$-امین مقدار منفرد مرتب‌شده ماتریس به صورت نزولی است. نسبت حفظ انرژی طیفی $\eta_r = \frac{\sum_{i=1}^r \sigma_i^2}{\sum_{i=1}^{d} \sigma_i^2}$ تضمین می‌کند که با انتخاب آستانه $\tau \ge 0.90$، بیش از نود درصد از پراکندگی واریانس داده‌ها حفظ شده و افت سیگنال در انتشار گرادیان رو به عقب به حداقل ممکن کاهش می‌یابد.

\paragraph{تحلیل مجانبی و پیچیدگی اشتراک حافظه سخت‌افزاری}
از منظر پیچیدگی محاسباتی و انتقال داده‌ها در سلسومراتب حافظه (حافظه سراسری \lr{HBM} به حافظه محلی \lr{SRAM})، عملیات ضرب ماتریسی در این مکانیزم نیازمند تخصیص بافر موقت با ابعاد فاکتورگیری‌شده است. برای پردازش دنباله‌ای به طول $N$ با بعد مخفی $D$، بار محاسباتی هسته از فرم درجه دو به رابطه زیر تقلیل می‌یابد:
\begin{equation}
\label{eq:complexity_sdrm_compressor}
\mathcal{C}_{\text{sdrm\_compressor}} = \mathcal{O}\left( B \cdot N \cdot D \cdot r + B \cdot H \cdot N \cdot d_k \right) \ll \mathcal{O}\left( B \cdot H \cdot N^2 \cdot d_k \right)
\end{equation}
این بهبود مرتبه پیچیدگی، با کاهش چشمگیر حجم تبادل تانسورها در گذرگاه‌های ارتباطی کارت گرافیک، مانع از اشباع پهنای باند حافظه (پدیده \lr{Memory-Bound}) در فرآیند استنتاج طولانی‌مدت می‌گردد.

\subsection{معماری و پیاده‌سازی}
پیاده‌سازی ماژول در فایل \lr{Codes/framework/compression\_engine/sdrm\_compressor.py} مستقر بوده و تحت نام‌های \lr{sdrm\_class} و \lr{sdrm\_compressor\_class} در رجیستری فریم‌ورک ثبت گردیده است. مشخصات فنی این پیاده‌سازی عبارتند از:
\begin{itemize}
    \item \textbf{مدیریت یکپارچه لایه‌های متوالی و ماژول‌های اشتراکی:} الگوریتم با بهره‌گیری از مجموعه \lr{visited\_modules} از پردازش تکراری لایه‌های با وزن مشترک (\lr{Tied Weights}) جلوگیری می‌کند.
    \item \textbf{کاهش فیزیکی واقعی در ابعاد لایه‌ها:} بر خلاف روش‌های تنک‌سازی که نیازمند پشتیبانی نرم‌افزاری خاص هستند، ساختار \lr{SVDLinear} در \lr{SDRM} ابعاد تانسورهای ذخیره‌شده را به طور فیزیکی کاهش می‌دهد.
    \item \textbf{وارسی پایداری و آستانه نرم اولیه:} پیش از تجزیه، نرم فروبنیوس ماتریس اولیه ارزیابی شده و لایه‌های با دامنه مقادیر نزدیک به صفر پردازش نمی‌شوند.
    \item \textbf{پشتیبانی کامل از لایه‌های کانوولوشن ۱ بعدی و خطی:} ابعاد ورودی و خروجی لایه‌های \lr{Conv1D} در مدل‌هایی نظیر \lr{GPT-2} با انطباق ترانهش ابعاد مدیریت می‌گردند.
\end{itemize}

\subsection{نتایج تجربی}
یافته‌های تجربی حاصل از بنچ‌مارک ماژول \lr{sdrm\_compressor} بر روی مدل‌های زبانی استاندارد در جدول \ref{tab:benchmark_sdrm} منعکس شده است:

\begin{table}[H]
\centering
\small
\caption{نتایج بنچ‌مارک تجربی ماژول فشرده‌ساز مدولاسیون رتبه \lr{sdrm\_compressor}}
\label{tab:benchmark_sdrm}
\begin{tabular}{lccccc}
\toprule
\textbf{مدل پایه \lr{(Base Model)}} & \textbf{دقت \lr{(Accuracy)}} & \textbf{شاخص \lr{F1}} & \textbf{حافظه \lr{(Peak VRAM)}} & \textbf{توان \lr{(Throughput)}} & \textbf{وضعیت} \\
\midrule
\lr{DistilBERT} & \lr{85.78\%} & \lr{0.8631} & \lr{338.98 MB} & \lr{32,743.48 tok/s} & \lr{PASS} \\
\lr{RoBERTa} & \lr{50.92\%} & \lr{0.6748} & \lr{781.79 MB} & \lr{16,739.56 tok/s} & \lr{PASS\_LOW} \\
\lr{GPT-2} & \lr{67.66\%} & \lr{0.7482} & \lr{446.90 MB} & \lr{10,921.04 tok/s} & \lr{PASS} \\
\lr{Qwen} & -- & -- & -- & -- & \lr{SKIP} \\
\bottomrule
\end{tabular}
\end{table}

بر پایه نتایج به دست آمده:
\begin{itemize}
    \item در مدل \lr{DistilBERT}، با فشرده‌سازی فیزیکی ۳۶ لایه و کاهش حجم پارامترها با نرخ فشرده‌سازی حقیقی \lr{1.24x}، دقت پایدار \lr{85.78\%} و شاخص \lr{F1} معادل \lr{0.8631} به دست آمده و وضعیت \lr{PASS} ثبت شده است.
    \item در مدل \lr{GPT-2}، با فشرده‌سازی ۱۵۲ لایه ساختاری و دستیابی به کاهش پارامتر با نسبت \lr{1.26x}، مدل دقت قابل توجه \lr{67.66\%} و شاخص \lr{F1} معادل \lr{0.7482} را در پیش‌بینی توکن‌ها کسب کرده است.
    \item در مدل \lr{RoBERTa}، به دلیل حساسیت شدید لایه‌های توجه این مدل نسبت به فشرده‌سازی فیزیکی ۶۸ لایه بدون فاین‌تیونینگ تکمیلی، دقت مدل به \lr{50.92\%} کاهش یافته که در سطح خطای تصادفی ارزیابی شده و وضعیت \lr{PASS\_LOW} را ایجاد نموده است.
    \item اجرای مدل \lr{Qwen} به علت محدودیت‌های زمانی ارزیابی در این نسخه انجام نگرفته و در وضعیت \lr{SKIP} گزارش شده است.
\end{itemize}

\subsection{محدودیت‌ها و تحلیل شکست}
تحلیل شکست و محدودیت‌های ماژول فشرده‌ساز \lr{SDRM} شامل موارد زیر است:
\begin{itemize}
    \item \textbf{افت دقت در مدل‌های حساس بدون گام بازیابی:} در معماری‌هایی مانند \lr{RoBERTa} که وزن‌ها همبستگی‌های غیرخطی عمیقی دارند، کاهش رتبه متناسب با بودجه در تمام لایه‌ها به صورت هم‌زمان موجب افت دقت به سطح تصادفی (\lr{PASS\_LOW}) می‌گردد که نشان می‌دهد برخی لایه‌های آغازین باید از فشرده‌سازی مستثنی شوند.
    \item \textbf{افت توان عملیاتی ناشی از سربار زنجیره‌ای ضرب ماتریسی دوگانه:} اگرچه تعداد محاسبات و پارامترها به میزان ۳۰ درصد کاهش یافته، جایگزینی یک ضرب ماتریسی متراکم بزرگ با دو ضرب ماتریسی متوالی در ابعاد کوچک می‌تواند در کارت‌های گرافیک سریع به دلیل سربار فراخوانی کرنل، گذردهی توکن را کاهش دهد.
\end{itemize}
